\documentclass{article}
\usepackage{amsmath,amssymb,hyperref}
\title{A stopped robot as both a clearance constraint and an optical occluder}
\date{26 November 2026}
\begin{document}
\maketitle
\paragraph{Claim and scope.}
In a civil inspection corridor a stopped robot can leave the required
one-relay bridge impossible at \emph{every} nominal range even when two
relays at a finite range suffice. In another regime each side has a valid
one-relay position, but switching sides under a requirement of continuous
two-terminal contact is impossible. A particular boundary family has an
exact range and speed limit for the switch when optical interruption is
permitted. These are model-level R6 physical feasibility checks, usable
before a coordination decision. They do not show a failed robot can be
localized, which direction an inspector should choose under partial
information, or that a visible link delivers a packet.

\paragraph{Collision and contact are different predicates.}
The terminal points are \(P=(-a,0)\) and \(Q=(a,0)\), \(a>0\).
The intervening stopped vehicle has an effective opaque disk with radius
\(0<r<a\). An inspecting relay center obeys clearance radius \(q\) with
\(r\le q<a\) and corridor center half-height \(H\ge0\).
The lateral dimension is unconstrained. These parameters already absorb
the stipulated obstacle, body and wall margins; they need not be equal.
An edge joins points at distance at most \(R\) if their segment does not
enter the open opaque disk. Tangency is allowed at the \emph{certified}
radius, for example if \(r\) is an inflated optical margin; strict
obstruction would change boundary cases. Anchors are optical points
without additional body-clearance constraints for the one-relay result.
Static placements presume known geometry and availability at that
position, not an authorized route to it. Optical orientations are
unconstrained. No map uncertainty, velocity matching at endpoints,
nonholonomic dynamics, interference, packet scheduling or accidental
physical optical access is assumed.

\paragraph{Exact one-relay feasibility.}
By reflection take location \((x,y)\) in the first quadrant. For
\(0\le t\le1\), the point at fraction \(t\) along the chord from \(P\) is
no farther from the origin than the corresponding point from \(Q\):
the squared distance difference is \(4tx(1-t)a\ge0\).
The segment from \(P\) is thus the harder visibility check and
also the longer link. The minimizing projection parameter onto its
support line is \(a(a+x)/((a+x)^2+y^2)\in[0,1]\).
The resulting squared distance to the segment is
\(a^2y^2/((a+x)^2+y^2)\). Write \(k=r^2/(a^2-r^2)>0\).
Both edges are optical exactly when \(y^2\ge k(a+x)^2\).

The feasible interval of \(z=y^2\) at fixed \(x\ge0\) is
\[
 [\max(q^2-x^2,k(a+x)^2),\ \min(H^2,R^2-(a+x)^2)] .
\]
The left endpoint is positive. Nonemptiness requires four pairwise
comparisons between the two lower and two upper bounds.
The clearance/height comparison says \(x\ge b=\sqrt{\max(0,q^2-H^2)}\).
Each of the remaining three gets weakly harder as \(x\) grows:
\[
 k(a+x)^2\le H^2,\qquad a^2+q^2+2ax\le R^2,\qquad
 (1+k)(a+x)^2\le R^2 .
\]
Consequently testing these at \(x=b\) is necessary and sufficient.
The left endpoint of the interval at \(b\), square-rooted, constructs
a valid \(y\), with reflected versions also valid. For rational inputs
the tests reduce to comparisons of \(\alpha+\beta\sqrt{B}\) with a
rational, where \(\beta\ge0\); check the residual sign before squaring.
When \(H\ge q\), the formula becomes \(H^2\ge ka^2\) and
\(R^2\ge a^2+\max(q^2,ka^2)\). The axial ray is optimal here, but
assuming axial placement when \(H<q\) can miss an actual bridge.

\paragraph{Separating witnesses.}
Take \(a=3,\ r=1,\ q=21/20,\ H=3/2\). At \(R^2=1011/100\),
a distance-only rule can approve \((0,q)\). The optical rule needs at
least \(R^2=81/8\). Now set \(H=211/200\), which is below the optical
height \(3/\sqrt 8\). No single optical relay works at any range.
The chain through \((-3/2,21/20)\) and \((3/2,21/20)\) works at \(R=3\):
every link has strictly positive optical margin, the relay centers
satisfy the body and wall clearance, and the four points are pairwise
at least unit distance apart. No assertion of optimal two-relay
\emph{range} or maintainable upstream deployment follows. The capacity
gain here comes from adding a relay, not making a stronger controller
at matched physical resources. Independently, for \(a=13,\ r=1,\ q=5,\ H=4\)
off-axis \((3,4)\) really works at \(R^2=272\); fixed axial placement
cannot obey clearance and height. At the same \(a,q,H\) but \(r=5\),
even the off-axis optimum fails the optical height check.

\paragraph{Static feasible endpoints do not give a live switch.}
Every allowed one-relay optical position has
\(|y|\ge h_0=a\sqrt{k}>0\), independently of \(R\).
Any continuous change from a positive to a negative position must pass
through \(y=0\), losing simultaneous visibility. If relay speed is
bounded by \(V>0\), a contiguous interruption between the last
positive-side service and the first negative-side service of at least
\(2h_0/V\) is necessary. This is a lower bound, not a universal optimum.
If both ends must remain within nominal distance during the movement,
crossing the axis outside \(q\) forces \(R\ge a+q\).
When \(H\ge\max(q,h_0)\), that and availability of the endpoints
\((0,\pm\max(q,h_0))\) suffice for those endpoints: travel along the
vertical rays and one semicircle of radius \(q\).

In the boundary family \(H=q=h_0<a\), the two service stations are the
\emph{only} feasible positions, so an outage budget \(W\) permits
such a switch, retaining both nominal distance bounds, exactly when
\(R\ge a+q\) and \(W\ge\pi q/V\). Indeed a rectifiable path outside the
clearance disk has length at least \(q\) times its net angular change
in absolute value, by the polar speed formula. Endpoint angles differ
by an odd multiple of \(\pi\); the right semicircle achieves the bound
at speed \(V\). All its interior points fail optical service.
For \(a=5,r=3\) this gives \(q=H=15/4\): radius \(25/4\) already serves
either fixed station, but radius \(35/4\) is needed for a switch under
the nominal distance rule. This is an explicitly tight geometric
limit, not a gap due to finding a sufficiently fine time step.
A spatial grid not containing \(x=0,y=\pm q\) can miss feasibility at
that exact boundary; a geometry-adaptive planner may insert the stations.

\paragraph{Prior mechanism, comparisons and residuals.}
\href{https://doi.org/10.1016/j.jfranklin.2023.05.013}{Maeda, Endo and Matsuno}
address local leader-following, a stopped robot, reconstruction and
collision avoidance.
\href{https://en.endo-lab.jp/research/detail/499}{The laboratory description}
identifies pre-stop direction and local waiting as ingredients of its
navigation work. The more recent
\href{https://doi.org/10.1109/TASE.2026.3741309}{Endo et al. distributed MPC}
also treats local sensing, collision avoidance and connectivity recovery.
Only scope from the published abstracts and lab descriptions is used;
no full controller reproduction or theorem subsumption is claimed.
Maintaining connectivity by explicitly testing segment visibility is
already used in
\href{https://experts.umn.edu/en/publications/visibility-based-deployment-of-robot-formations-for-communication/}{visibility-based deployment}.
A direct exact planner given the same continuous constraints cannot beat
the infeasibility result and can use its constructive position witness.
This is a resource and maneuver boundary for one restricted instance
family, not an advantage over such a planner. The older project
\href{https://github.com/eternal-november/eos-research-programmes/pull/396}{relay lens}
used separate range and obstruction of motion; here the same obstacle
removes segment visibility in otherwise navigable configurations.

Lean checks the chord comparison, algebraic visibility test and two
elimination steps. Exact rational ray/segment calculations check
finite point, threshold, slot and trajectory sample controls. The
analytical argument, not the samples, establishes continuum feasibility
and the speed bound. Installation timing, active observation, report
custody and acknowledgment, retained route clearance, and whether an
independent failover source can provide an equivalent acceptance
interface remain separate. No inventory or contract disposition changes.
\end{document}
